\chapter{Lua Scripting}

OmniLaTeX embeds Lua\TeX{} scripting to provide dynamic metadata,
conditional compilation, and build-time automation.

\section{Lua in \LaTeX}

Lua code can be executed inline with \texttt{\textbackslash directlua} or
in a dedicated environment:

\begin{verbatim}
\directlua{tex.sprint("2 + 2 = ", 2 + 2)}

\begin{lua}
    local function greeting(name)
        return "Hello, " .. name .. "!"
    end
    tex.sprint(greeting("World"))
\end{lua}
\end{verbatim}

\texttt{\textbackslash luasetup} is provided for one-time initialisation
blocks that run before the document body.

\section{Git Metadata}

\texttt{omnilatex-utils.sty} queries the working tree at compile time and
exposes the following macros:

\begin{table}[htbp]
    \centering
    \caption{Git metadata macros}
    \label{tab:git-macros}
    \begin{tabular}{@{} l p{7cm} @{}}
        \toprule
        \textbf{Macro} & \textbf{Expansion} \\
        \midrule
        \verb|\GitRefName|   & Current branch or tag name \\
        \verb|\GitShortSHA|  & Abbreviated commit hash (8 hex digits) \\
        \verb|\GitLongSHA|   & Full 40-character SHA \\
        \verb|\BuildDate|    & ISO-8601 build timestamp \\
        \verb|\BuildMode|    & String: \texttt{dev}, \texttt{prod}, or
                               \texttt{ultra} \\
        \bottomrule
    \end{tabular}
\end{table}

These macros expand to empty strings when Git metadata cannot be resolved
(e.g.\ in CI environments without a \texttt{.git} directory).

\section{Conditional Compilation}

The build mode is determined by the \texttt{BUILD\_MODE} environment
variable.  OmniLaTeX defines three boolean flags:

\begin{verbatim}
\newif\ifomnilatex@dev
\newif\ifomnilatex@prod
\newif\ifomnilatex@ultra
\end{verbatim}

Detection is performed at load time and exposed through user-facing
switches:

\begin{verbatim}
\OmniDev   % active when BUILD_MODE=dev
\OmniProd  % active when BUILD_MODE=prod
\OmniUltra % active when BUILD_MODE=ultra
\end{verbatim}

Use these in the document body to include or exclude content:

\begin{verbatim}
\ifomnilatex@dev
    \typeout{[dev] Draft watermark enabled.}
\fi
\end{verbatim}

\section{Custom Lua Commands}

Arbitrary Lua functions can be wrapped as \LaTeX{} commands with the
\texttt{\textbackslash directlua\textbackslash stringdef} pattern:

\begin{verbatim}
\directlua\stringdef\mycmd#1{%
    tex.sprint(string.upper(#1))%
}
\end{verbatim}

\texttt{\textbackslash mycmd\{hello\}} then expands to \texttt{HELLO}.

\section{Example Walkthrough}

A complete demonstration of Lua scripting, including metadata queries,
conditional blocks, and custom commands, is available in the
\texttt{examples/lua-showcase/} directory.  Compile that sub-project with

\begin{verbatim}
latexmk -lualatex examples/lua-showcase/main.tex
\end{verbatim}
